% ============================================================
%  SUPPLEMENT 01 — ECDR PRIMER
%  What ECDR Is, How It Works, and Why It Is Sound
%  Companion to: ECDR_Data_Analysis.tex
%  Do NOT modify ECDR_Data_Analysis.tex — add only here.
% ============================================================
\documentclass[11pt, a4paper]{article}

% ---------- fonts & encoding ----------
\usepackage{fontspec}
\setmainfont{Noto Serif}
\setsansfont{Noto Sans}
\setmonofont{Noto Sans Mono}[Scale=0.88]

% ---------- language ----------
\usepackage[english]{babel}

% ---------- math ----------
\usepackage{amsmath, amssymb, amsthm}
\usepackage{mathtools}

% ---------- layout ----------
\usepackage[margin=2.5cm]{geometry}
\usepackage{booktabs, tabularx, array, longtable}
\usepackage{xcolor}
\newcolumntype{Y}{>{\raggedright\arraybackslash}X}
\newcolumntype{P}[1]{>{\raggedright\arraybackslash}p{#1}}
\usepackage{hyperref}
\usepackage{enumitem}
\usepackage{microtype}
\usepackage{parskip}

% ---------- color scheme (matches main doc) ----------
\definecolor{ink}{RGB}{20,20,20}
\definecolor{teal}{RGB}{0,120,120}
\definecolor{muted}{RGB}{90,90,90}
\definecolor{lightgray}{RGB}{240,240,240}

\hypersetup{colorlinks=true, linkcolor=teal, urlcolor=teal, citecolor=teal}

% ---------- theorem environments ----------
\theoremstyle{definition}
\newtheorem{definition}{Definition}[section]
\newtheorem{example}{Example}[section]
\newtheorem{remark}{Remark}[section]

\theoremstyle{plain}
\newtheorem{theorem}{Theorem}[section]
\newtheorem{lemma}{Lemma}[section]

% ---------- macros ----------
\newcommand{\R}{\mathbb{R}}
\newcommand{\E}{\mathbb{E}}
\newcommand{\Prob}{\mathbb{P}}
\newcommand{\Se}{\mathcal{S}_e}
\newcommand{\code}[1]{\texttt{#1}}

% ---------- section formatting ----------
\usepackage{titlesec}
\titleformat{\section}{\large\bfseries\color{teal}}{}{0em}{\thesection\quad}
\titleformat{\subsection}{\normalsize\bfseries\color{ink}}{}{0em}{\thesubsection\quad}

% ============================================================
\begin{document}

\begin{center}
  {\LARGE\bfseries\color{teal} ECDR Primer}\\[0.5em]
  {\large\color{muted} What ECDR Is, How It Works, and Why It Requires Formal Certification}\\[1.2em]
  {\normalsize Supplement~01 to \textit{ECDR Data Analysis} \textbar{} ECDR Research Workspace, October~2026}\\[0.3em]
  {\small\color{muted}
    This document is a supplementary explanation only. \\
    The main data analysis document (\code{ECDR\_Data\_Analysis.tex}) is not modified.}
\end{center}

\vspace{1em}
\hrule
\vspace{1em}

\tableofcontents
\clearpage

% ============================================================
\section{Motivation: Why Reuse Decisions at All?}

Neural networks deployed on edge devices (smartphones, wearables, embedded sensors)
face two contradictory requirements:

\begin{enumerate}
  \item \textbf{Accuracy}: The activity recognition model must classify each incoming
        sensor window correctly, which demands running the full neural network inference
        pipeline on every query.
  \item \textbf{Efficiency}: Edge hardware has severe power and latency budgets.
        Running a full CNN forward pass at 50\,Hz drains battery and heats the device.
\end{enumerate}

The key insight behind ECDR is empirical: \emph{when a person walks steadily, consecutive
1.28\,s windows look almost identical to the network.} The network would produce the same
logits, and therefore the same top-1 decision, on many consecutive queries. If we can
\emph{prove} that the decision cannot change within a neighborhood, we can safely skip the
forward pass for any query that falls inside that neighborhood.

\begin{remark}[ECDR Is Not a Cache]
  A software cache stores results keyed by exact input hashes: only when
  $X_t = X^r$ byte-for-byte does a cache hit occur. ECDR is a \emph{certified
  neighborhood reuse} system: it guarantees that the classification
  decision is constant over a \emph{continuous, infinite-dimensional neighborhood}
  around $X^r$. This covers infinitely many possible sensor readings, not just
  one exact recording.
\end{remark}

% ============================================================
\section{The Core Invariance Argument}

\subsection{Setup and Notation}

Let the sensor stream produce a sequence of windows
$X_1, X_2, \ldots \in \R^{6 \times 128}$, where each window contains 128 time
steps of 6-channel IMU data (3-axis accelerometer and 3-axis gyroscope).

The full inference pipeline consists of two stages:
\begin{enumerate}
  \item \textbf{Normalization}: $N: \R^{6\times128} \to \R^{6\times128}$ is a
        fixed, affine channel-wise standardization (subtract mean, divide by
        standard deviation). It is determined once from training data and frozen.
  \item \textbf{Neural network}: $f_\theta: \R^{6\times128} \to \R^K$ maps a
        normalized window to raw logits over $K$ activity classes.
\end{enumerate}

The composed classifier is $F \coloneqq f_\theta \circ N$. The top-1 decision is
$c(X) \coloneqq \arg\max_{j \in \{1,\ldots,K\}} F_j(X)$.

\subsection{What Is an ``Anchor''?}

An \textbf{anchor} is a specific window $X^r$ that has been fully evaluated by the
reference classifier $F$, yielding top-1 decision $c = c(X^r)$ and the full logit
vector $F(X^r) \in \R^K$.

During streaming, every subsequent window $X_t$ for $t > r$ is a \textbf{query}.
Rather than running $F$ on $X_t$, ECDR checks: \emph{can we prove that
$c(X_t) = c$?} If yes, reuse the decision. If not (or if the certificate expires),
evaluate $F(X_t)$ directly and potentially set a new anchor.

\subsection{The Certified Reuse Condition}

\begin{definition}[Perturbation Domain]
  A perturbation domain $\Se(X^r)$ is a measurable subset of $\R^{6\times128}$
  centered at the anchor $X^r$, designed to contain plausible streaming queries
  that are ``close'' to $X^r$ in a physically meaningful sense.
\end{definition}

\begin{definition}[Logit Margin Lower Bound]
  For a given domain $\mathcal{S}$ and class pair $(c, j)$ with $j \ne c$, define
  the \textbf{certified margin}:
  \begin{equation}
    \underline{m}_{cj}(\mathcal{S})
    \;\le\; \inf_{X \in \mathcal{S}} \bigl[F_c(X) - F_j(X)\bigr].
  \end{equation}
  The quantity $\underline{m}_{cj}$ is a \emph{sound lower bound}: it is guaranteed
  to be at most the true minimum logit gap between class $c$ and class $j$ over the
  entire domain. If $\underline{m}_{cj} > 0$, then $F_c(X) > F_j(X)$ for every
  $X \in \mathcal{S}$, meaning class $c$ always wins over class $j$ inside $\mathcal{S}$.
\end{definition}

\begin{theorem}[Sound Decision Reuse]
  If\,
  $X_t \in \Se(X^r)$
  \;\textbf{and}\;
  $\underline{m}_{cj}\bigl(\Se(X^r)\bigr) > 0$ for all $j \ne c$,
  \;then $c(X_t) = c$.
\end{theorem}

\begin{proof}
  For any $X_t \in \Se(X^r)$ and any competing class $j \ne c$:
  \[
    F_c(X_t) - F_j(X_t)
    \;\ge\; \inf_{X \in \Se(X^r)}\bigl[F_c(X) - F_j(X)\bigr]
    \;\ge\; \underline{m}_{cj} > 0.
  \]
  Therefore $F_c(X_t) > F_j(X_t)$ for all $j \ne c$, so $\arg\max_j F_j(X_t) = c$.
\end{proof}

\begin{remark}[Both Conditions Are Necessary]
  \begin{itemize}[leftmargin=1.5em]
    \item $X_t \in \Se(X^r)$ alone: we know $X_t$ is nearby, but the network may
          still flip its decision inside $\Se$ — no safety guarantee.
    \item $\underline{m}_{cj} > 0$ alone: the margin is certified over the domain,
          but we have not verified that the specific query $X_t$ falls inside — no
          applicability guarantee.
    \item \textbf{Both together}: $X_t$ is inside the certified region, and the
          decision is proven constant throughout that region. Safe reuse.
  \end{itemize}
\end{remark}

% ============================================================
\section{ECDR vs.\ Simple Caching: A Precise Comparison}

\begin{center}
\small
\begin{tabularx}{\linewidth}{@{}lYY@{}}
\toprule
Property & Hash Cache & ECDR Certified Reuse \\
\midrule
Hit condition & $X_t = X^r$ (exact bitwise match) & $X_t \in \Se(X^r)$ (continuous neighborhood) \\
Hit probability & Near zero for continuous sensors & Tunable via domain size \\
Safety guarantee & Trivially exact & Formally proven by verifier \\
Domain shape & Point (no neighborhood) & Hyperrectangle, $L_2$ ball, wave-domain, etc. \\
Requires verifier? & No & Yes (IBP/CROWN/LiRPA) \\
Effective reuse & Negligible in streaming & Potentially $>90\%$ for static postures \\
\bottomrule
\end{tabularx}
\end{center}

% ============================================================
\section{Anchor Lifecycle: States and Transitions}

An anchor does not live forever. As time progresses, the streaming query $X_t$ may
drift away from $X^r$, eventually leaving $\Se(X^r)$. The system must track the
status of each anchor and know when to retire it.

\subsection{State Machine}

\begin{center}
\begin{tabular}{@{}lll@{}}
\toprule
State & Meaning & Allowed Actions \\
\midrule
\textbf{FRESH} & Certificate just constructed on $X^r$; margin verified. & Reuse if $X_t \in \Se$. \\
\textbf{ACTIVE} & Anchor is in use; successive queries accepted. & Reuse if $X_t \in \Se$. \\
\textbf{STALE} & Anchor has been active for too long, or drift is detected. & Do not reuse. Must refresh. \\
\textbf{REFRESH} & Running full $F(X_t)$ and re-certifying. & Blocking computation step. \\
\midrule
\multicolumn{3}{@{}l}{\textit{Transitions:} FRESH $\to$ ACTIVE (first reuse), ACTIVE $\to$ STALE (timeout / domain exit),} \\
\multicolumn{3}{@{}l}{\phantom{\textit{Transitions:}} STALE $\to$ REFRESH $\to$ FRESH (new anchor)} \\
\bottomrule
\end{tabular}
\end{center}

\subsection{Why Recentering Is Forbidden Without Re-Verification}

An important failure mode: \emph{slide the anchor center to $X_t$ after each accepted
query.}

Suppose $X_t$ was accepted inside $\Se(X^r)$. Now redefine $X^r \leftarrow X_t$
and build a new domain $\Se(X_t)$ \emph{without re-running the verifier}. This new
domain extends in directions outside the original $\Se(X^r)$. The logit margin over
those new regions has never been certified. A competing class might cross above $F_c$
just outside $\Se(X^r)$ but inside $\Se(X_t)$.

\begin{remark}[Formal Statement]
  $\Se(X_t)$ need not be a subset of $\Se(X^r)$ even when $X_t \in \Se(X^r)$.
  Certification of $\Se(X^r)$ implies nothing about $\Se(X_t)$.
  Re-verification is a non-negotiable requirement.
\end{remark}

\subsection{Drift and the Triangle Inequality}

For a sequence of queries $X_{r}, X_{r+1}, \ldots, X_{r+h}$ (horizon $h$ steps),
the accumulated drift satisfies:
\begin{equation}
  \|X_{r+h} - X^r\|
  \;\le\; \sum_{k=0}^{h-1} \|X_{r+k+1} - X_{r+k}\|.
\end{equation}
Each one-step $L_2$ distance is bounded empirically (median $\approx 4.9$ normalized
for same-activity pairs; median $\approx 36.6$ for locomotion classes). For static
postures the drift is negligible; for walking, $h = 5$ steps can accumulate
$\approx 5 \times 4.9 = 24.5$ normalized units of displacement — enough to exit
typical certification radii.

This means \textbf{anchor lifetime depends strongly on activity class}, a key
empirical finding of the data analysis.

% ============================================================
\section{The Two Branches of the HAPT Feature Space}

The HAPT dataset is processed in two distinct formats within the ECDR workspace:

\subsection{Raw 6-Channel Windows (This Document)}
\begin{itemize}[leftmargin=1.5em]
  \item Shape: $6 \times 128$ tensor (6 IMU channels, 128 time steps at 50\,Hz).
  \item The perturbation domain $\Se$ lives in $\R^{768}$.
  \item ECDR certification operates directly on continuous sensor space.
  \item Primary focus of \code{ECDR\_Data\_Analysis.tex}.
\end{itemize}

\subsection{561-Feature Engineered Branch}
\begin{itemize}[leftmargin=1.5em]
  \item The original HAPT data release also provides a 561-dimensional feature
        vector per window, computed by the dataset authors via fixed-window FFT,
        energy bands, gravity component separation, Jerk signals, and other
        hand-crafted transformations.
  \item A classifier operating on 561-features occupies $\R^{561}$ — a completely
        different space with different geometry, different norm relationships,
        and different certification requirements.
  \item The two branches are \textbf{not interchangeable}: empirical distances
        measured in $\R^{768}$ raw space do not translate to $\R^{561}$ feature
        space, and vice versa.
  \item ECDR certification in the 561-feature branch requires a separate
        normalization pipeline, separate verifier invocation, and separate
        domain construction.
\end{itemize}

% ============================================================
\section{What Makes Certification ``Sound''?}

The word \emph{sound} has a precise technical meaning in formal verification:

\begin{definition}[Soundness]
  A procedure for computing $\underline{m}_{cj}(\mathcal{S})$ is \textbf{sound}
  if the computed value is guaranteed to be a lower bound on the true infimum:
  \[
    \underline{m}_{cj}(\mathcal{S}) \;\le\; \inf_{X \in \mathcal{S}} \bigl[F_c(X) - F_j(X)\bigr].
  \]
  A sound procedure may under-estimate the margin (producing a smaller value than
  the truth) but will never over-estimate it (claiming the margin is larger than
  it actually is).
\end{definition}

\begin{itemize}[leftmargin=1.5em]
  \item \textbf{Why this matters:} If a verifier accidentally reports
        $\underline{m}_{cj} > 0$ when the true infimum is $\le 0$, the system
        would reuse a decision that is actually wrong on some $X \in \mathcal{S}$.
        That is a \emph{safety failure}.
  \item \textbf{Tightness is efficiency:} A sound verifier that always outputs
        $\underline{m}_{cj} = -\infty$ is trivially sound but useless — it
        certifies nothing. Good verifiers (CROWN, LiRPA) produce tight lower
        bounds, certifying large domains.
  \item \textbf{Completeness is different:} A \emph{complete} verifier returns
        $\underline{m}_{cj} > 0$ if and only if the true margin is positive.
        Neural network verification is in general NP-hard, so practical tools
        sacrifice completeness for tractability while preserving soundness.
\end{itemize}

\subsection{Cryptographic Hash Checks Are Not Certification}

The data audit uses SHA-256 hashes to verify that tensor files have not been
accidentally corrupted. This is a \emph{data integrity} check, not a
\emph{semantic correctness} guarantee. A SHA-256 match confirms that two files
are bitwise identical; it says nothing about:
\begin{itemize}[leftmargin=1.5em]
  \item Whether the annotation indexing scheme (0-based vs.\ 1-based) is correct.
  \item Whether the normalization parameters are appropriate.
  \item Whether the classifier's margins are certified over $\Se$.
\end{itemize}
Cryptographic verification and formal neural network certification are
complementary but entirely distinct mechanisms.

% ============================================================
\section{The ECDR Safety Argument: End-to-End Summary}

The full safety argument for a single reuse event is:

\begin{center}
\small
\begin{tabular}{@{}clll@{}}
\toprule
Step & Who Performs It & What Is Established & Evidence Type \\
\midrule
1 & Data pipeline & $X_t$ correctly extracted from raw IMU & SHA-256 + float32 cast check \\
2 & Normalization $N$ & $Z_t = N(X_t)$ correctly computed & Arithmetic: affine, closed-form \\
3 & Membership test & $X_t \in \Se(X^r)$ checked & Runtime coordinate comparison \\
4 & Formal verifier & $\underline{m}_{cj}(\Se) > 0$ for all $j \ne c$ & Sound IBP/CROWN bound \\
5 & System logic & Reuse $c$ as output & Follows from Theorem~1 \\
\midrule
\multicolumn{4}{@{}l}{If any step fails: fall back to full inference $F(X_t)$.} \\
\bottomrule
\end{tabular}
\end{center}

Breaking any one step invalidates the guarantee. The data audit
(\code{ECDR\_Data\_Analysis.tex}) provides empirical support for Steps~1 and~3.
Steps~4 and~5 require separate verifier implementation and are currently at
\textbf{HYPOTHESIS} epistemic status (D07, D08 in the Evidence Ledger).

% ============================================================
\section{Open Problems as of This Document}

\begin{longtable}{@{}llp{7cm}@{}}
\toprule
Ref & Status & Description \\
\midrule
D06 & OPEN & Annotation indexing (0-based vs. 1-based): upstream intent unresolved. \\
D07 & HYPOTHESIS & Structured domains yield practical certified coverage for a 12-class CNN. \\
D08 & HYPOTHESIS & ECDR delivers net savings on edge hardware ($T_e + T_m + qT_b < rT_f$). \\
D09 & OPEN & Reference CNN maintains high accuracy across all 12 classes on HAPT. \\
\bottomrule
\end{longtable}

\noindent Until D07 and D08 are tested with an actual trained network and profiled
on target hardware, ECDR should be characterized as a \emph{theoretically sound,
empirically unvalidated} system for the 12-class setting.

\begin{thebibliography}{9}
\bibitem{main}
  ECDR Research Workspace. \textit{ECDR Data Analysis} (English edition).
  \code{ECDR\_Data\_Analysis.tex}, October~2026.
\bibitem{weeks}
  ECDR Research Workspace. \textit{ECDR Weeks~1--3 Theory Brief}.
  \code{ECDR\_Weeks\_1\_2\_3\_Theory.tex}, October~2026.
\bibitem{math}
  ECDR Research Workspace. \textit{ECDR Mathematics Concepts}.
  \code{ECDR\_Mathematics\_Concepts\_EN.tex}, October~2026.
\end{thebibliography}

\end{document}
